\documentclass[11pt]{article}
\usepackage[margin=1.1in]{geometry}
\usepackage{amsmath,amssymb,amsthm}
\usepackage{booktabs}
\usepackage[hidelinks]{hyperref}
\usepackage{enumitem}

\newtheorem{lemma}{Lemma}
\newtheorem{proposition}{Proposition}
\newtheorem{theorem}{Theorem}
\newtheorem{corollary}{Corollary}
\newtheorem{assumption}{Assumption}
\theoremstyle{remark}
\newtheorem{remark}{Remark}

\newcommand{\E}{\mathbb{E}}
\newcommand{\Pro}{\Pr}
\newcommand{\Cov}{\mathrm{Cov}}
\newcommand{\ind}[1]{\mathbf{1}\{#1\}}
\newcommand{\statusP}{\textsc{[proved]}}
\newcommand{\statusPA}{\textsc{[proved under assumptions]}}
\newcommand{\statusS}{\textsc{[sketch]}}

\title{Rebuilding the Deliberation Model from Primitives:\\
Derivations, Proofs, and an Honest Audit\thanks{Internal theory note for
Hall \& Huang, ``The Effect of AI Writing on Governance: Evidence from a
\$3.5 Billion DAO.'' Prepared for the Thursday meeting (Decision~1 of the
post-review agenda). Builds on V.~Huang's memos \texttt{model\_derivation\_notes.md}
and \texttt{model\_proofs.md}; all results here are re-derived from scratch and
verified numerically (script \texttt{check\_model.py}). Every result is labeled
\textsc{proved}, \textsc{proved under assumptions}, or \textsc{sketch}.}}
\author{Prepared for J.~Hall and V.~Huang}
\date{August 2026}

\begin{document}
\maketitle

\begin{abstract}
\noindent The referee's central objection is that the paper's model asserts rather
than derives its three headline objects: interior vote margins, posting volume
proportional to $p(1-p)$, and the welfare kernel $4p(1-p)(q-\tfrac12)^2$. This
note rebuilds the model as a unified two-tier structure --- informed contributors
post; a continuum of taste-heterogeneous voters decides whether to read and then
votes --- in which \emph{all} equilibrium objects (readership, posting volume,
vote margins, decision quality, and the AI tipping point) are derived from one
set of primitives. The main findings: (i)~Proposition~1 of the paper
($\Cov(\text{posts},\text{margin})<0$) is now a theorem; (ii)~the
career-concerns posting motive in the fix roadmap is \emph{provably wrong-signed}
for any increasing reputation payoff and must be replaced by an
attention/value-of-information motive; (iii)~the paper's reading rule
$V(\pi)=(1-\pi)I_H-c_r$ is derived \emph{exactly}, with
$I_H=2p(1-p)(2q-1)$; (iv)~the welfare formula $4p(1-p)(q-\tfrac12)^2$ does
\emph{not} survive --- the exact value of the forum is
$\min(p,1-p)$ on an interior ``flip window,'' zero outside; (v)~the tipping-point
result survives and is \emph{strengthened}: with endogenous contamination, a
positive reading cost alone (no S-shaped adoption curve needed) makes collapse
discontinuous, hysteretic, and located at a realized AI share strictly below the
reading threshold; (vi)~of the calibration, the bound $\pi^*\ge 0.082$ and the
19.4 words-per-post surplus survive; the $\times 44.2$ aggregation to 858
words/delegate/proposal survives only under an explicitly stated
non-redundancy assumption. A verdict for the Thursday meeting is in
Section~\ref{sec:verdict}.
\end{abstract}

\section{Purpose and the three objections}
\label{sec:intro}

The automated review makes three related objections to Section~5 of the paper:
\begin{enumerate}[label=(O\arabic*),itemsep=1pt]
\item a single representative delegate making a binary vote cannot generate an
interior, continuously varying margin of victory;
\item with a nonatomic continuum of contributors no one is pivotal, so costly
posting is not an equilibrium and ``posting volume'' has no microfoundation;
\item the welfare formula $W^*-W^{AI}=4p(1-p)(q-\tfrac12)^2$ is asserted, not
derived.
\end{enumerate}
All three share one root: the paper's model has one decision-maker where the
data have populations. The rebuild below makes both sides populations. A
continuum of privately informed \emph{contributors} decides whether to post; a
continuum of \emph{voters} with heterogeneous private tastes decides whether to
\emph{read} and then votes. Margins are interior because voters disagree about
values, not only facts (answering O1); posting is an equilibrium because posters
are paid in \emph{audience}, which is a positive mass even for a measure-zero
poster (answering O2, with the assumption doing the work stated explicitly); and
welfare is computed exactly from the equilibrium voting behavior (answering O3
--- with the consequence that the old formula must be replaced).

Two candidate posting motives were on the table in the fix roadmap. We prove
below (Theorem~\ref{thm:career}) that the reputation/career-concerns motive
generates posting volume \emph{minimized} on contested proposals --- the opposite
of the data --- for \emph{any} increasing reputation payoff, and therefore cannot
underlie Proposition~1. The attention/value-of-information motive works and is
adopted.

\section{The model}
\label{sec:model}

\subsection{Primitives}

A DAO evaluates one proposal. Its unknown quality is $\theta\in\{G,B\}$ with
common prior $p\equiv\Pro(\theta=G)\in(0,1)$. Write $x\equiv p(1-p)$ and
$d\equiv|p-\tfrac12|$; $d$ measures ex-ante \emph{consensus} ($d$ large) versus
\emph{contestedness} ($d$ small).

\textbf{Contributors.} A continuum of mass $1$ of potential posters. Contributor
$i$ observes a private signal $s_i\in\{g,b\}$ with symmetric precision
$q\equiv\Pro(s_i=g\mid G)=\Pro(s_i=b\mid B)\in(\tfrac12,1)$, i.i.d.\ conditional
on $\theta$, and draws a posting cost $c_i\sim F$ on $[0,\bar c]$, independent of
$s_i$, where $F$ is continuous and strictly increasing with $F(0)=0$. Posting
publishes a message that reveals $s_i$.

\textbf{AI content.} A mass $z\ge 0$ of AI-generated posts is present in the
thread; an AI post's message is uninformative
($\Pro(\text{message}=g\mid\theta)=\tfrac12$ for both $\theta$) and is
observationally indistinguishable from a human post. If human posting mass is
$n$, the fraction of AI posts is $\pi=z/(z+n)$. Sections
\ref{sec:read}--\ref{sec:welfare} treat $\pi$ as a parameter;
Section~\ref{sec:tipping} makes it an equilibrium object.

\textbf{Voters.} A continuum of mass $1$. Voter $j$ has a private taste
$b_j\sim\mathrm{Unif}[-a,a]$, i.i.d., orthogonal to $\theta$ (idiosyncratic
value from the proposal passing: exposure, vested interests, ideology). Common
values are $v(G)=+1$, $v(B)=-1$: voting \emph{For} a proposal that passes yields
$v(\theta)+b_j$, \emph{Against} yields $0$. Before voting, a voter may pay
$c_r\ge 0$ to read one post sampled uniformly at random from the thread
(Remark~\ref{rem:kpost} discusses multi-post reading).

\textbf{Timing.} (1) Nature draws $\theta$, signals, costs, tastes. (2)
Contributors post or not; AI mass $z$ is posted. (3) Each voter chooses to read
or not; readers sample one post each, independently. (4) All voters vote For or
Against; the proposal passes iff the For-share $\Phi$ exceeds $\tfrac12$.
(5) Payoffs realize.

\subsection{Assumptions, stated honestly}

\begin{assumption}[Exact LLN]\label{ass:lln}
Conditional on $\theta$, cross-sectional averages of i.i.d.\ individual
quantities equal their conditional means (the usual continuum convention; see
Judd 1985 for the measurability caveat). All ``deterministic'' aggregates below
are exact in this convention and hold up to $O(N^{-1/2})$ noise with $N$ finite
agents.
\end{assumption}

\begin{assumption}[Decision-theoretic voting and reading]\label{ass:sincere}
Each voter votes --- and values pre-vote information --- \emph{as if her vote
determined the outcome for her own payoff}: she votes For iff
$\E[v(\theta)\mid\text{her information}]+b_j\ge 0$, and values reading by the
improvement in this decision.
\end{assumption}

This is the standard sincere/expressive-voting device for large elections
(e.g., ethical-voter models in the spirit of Feddersen and Sandroni 2006). It
must be said plainly: \emph{the pivotality problem on the voter side is bypassed
by this assumption, not solved}. What the assumption buys is discipline: given
it, everything else (who reads, how vote shares move, when the forum dies) is
derived. The alternative --- fully instrumental voters --- makes reading and
voting trivially worthless for every atomless agent and cannot generate any of
the observed cross-sectional structure.

\begin{assumption}[Truthful posting]\label{ass:truth}
A post reveals the contributor's signal truthfully.
\end{assumption}

Justifiable by reputational cheap-talk arguments (posts are public and
$\theta$ is eventually revealed; systematic misreporting is detectable ex post,
cf.\ Ottaviani and S{\o}rensen 2006), but in this note it is an assumption.

\begin{assumption}[Attention rewards]\label{ass:attention}
A contributor who posts receives $R\cdot A$, where $A$ is the mass of readers of
the thread and $R>0$ converts audience into utility.
\end{assumption}

This is the posting motive that survives the continuum: audience is a positive
\emph{mass} even though any single poster is measure zero, so no pivotality is
needed. Microfoundations consistent with the setting: delegate incentive
programs that compensate visible participation; attracting delegated voting
power from readers; direct attention utility. What Assumption~\ref{ass:attention}
deliberately does \emph{not} assume is that posting signals ability ---
Theorem~\ref{thm:career} shows that motive predicts the wrong sign.
Remark~\ref{rem:congestion} shows the results are robust to rewarding
\emph{per-post} rather than per-thread attention.

\begin{assumption}[Interior tastes]\label{ass:interior}
$a\ge 1$, so every posterior-based voting threshold lies inside $[-a,a]$ and
vote shares are interior.
\end{assumption}

\begin{assumption}[Pre-AI viability]\label{ass:viable}
$c_r< I_H(p)\equiv 2p(1-p)(2q-1)$ for the proposals under study: absent AI,
reading has positive net value for at least the most contested proposals.
\end{assumption}

\subsection{Notation and two identities}

A sampled post is human with probability $1-\pi$; hence its message has
effective precision
\begin{equation}\label{eq:Qeff}
Q \;\equiv\; \Pro(\text{message}=g\mid G)\;=\;(1-\pi)\,q+\pi\cdot\tfrac12,
\qquad\text{so}\qquad Q-\tfrac12=(1-\pi)\bigl(q-\tfrac12\bigr).
\end{equation}
This is the ``signal dilution'' equation of the working notes, here \emph{derived}
from the sampling protocol rather than posited. Write
\[
D_1=\Pro(g)=pQ+(1-p)(1-Q),\quad D_2=\Pro(b)=p(1-Q)+(1-p)Q,\quad
\mu_G=\tfrac{pQ}{D_1},\quad \mu_B=\tfrac{p(1-Q)}{D_2},
\]
for the message probabilities and posteriors ($\mu_m=\Pro(G\mid m)$), and
$D\equiv D_1D_2$.

\begin{lemma}[Bayesian identities \statusP]\label{lem:identities}
\begin{align}
\textup{(I1)}\quad & D_1(\mu_G-p)=D_2(p-\mu_B)=p(1-p)(2Q-1),\\
\textup{(I2)}\quad & \mu_G-\mu_B=\frac{p(1-p)(2Q-1)}{D_1D_2},\\
\textup{(I3)}\quad & D_1D_2=Q(1-Q)+p(1-p)(2Q-1)^2.
\end{align}
\end{lemma}

\begin{proof}
(I1): $D_1\mu_G=pQ$, so $D_1(\mu_G-p)=pQ-pD_1=p[(1-p)Q-(1-p)(1-Q)]=x(2Q-1)$;
symmetrically for $D_2$. (I2): $\mu_G-\mu_B=(\mu_G-p)+(p-\mu_B)
=x(2Q-1)(D_1^{-1}+D_2^{-1})$ and $D_1+D_2=1$. (I3): expand
$D_1D_2=Q(1-Q)[p^2+(1-p)^2]+x[Q^2+(1-Q)^2]=Q(1-Q)(1-2x)+x[1-2Q(1-Q)]$ and
collect terms.
\end{proof}

Note $D\in(0,\tfrac14]$, with $D=\tfrac14$ at $p=\tfrac12$.

\section{The reading stage: who reads, and the paper's reading rule derived}
\label{sec:read}

Voters differ in how useful the forum is to them. A voter who does not read
votes on the prior: For iff $b_j\ge t_0\equiv 1-2p$. A reader who sees message
$m$ votes For iff $b_j\ge t_m\equiv 1-2\mu_m$; note $t_G<t_0<t_B$.

\begin{lemma}[Value of reading is a tent \statusP]\label{lem:tent}
The value of information from reading one post for a voter with taste $b$ is
\[
\mathrm{VOI}(b)=
\begin{cases}
D_1\,(b-t_G), & b\in[t_G,t_0],\\[2pt]
D_2\,(t_B-b), & b\in[t_0,t_B],\\[2pt]
0,& \text{otherwise:}
\end{cases}
\]
a tent on the swing band $[t_G,t_B]$ (base $2(\mu_G-\mu_B)$), with peak at the
ex-ante indifferent voter $b=t_0$ of height
\begin{equation}\label{eq:peak}
h(p,Q)\;=\;2p(1-p)(2Q-1).
\end{equation}
\end{lemma}

\begin{proof}
Reading changes the action only for $b\in[t_G,t_B)$. For $b\in[t_G,t_0)$ the
voter switches to For only on $m=g$, gaining $D_1[(2\mu_G-1)+b]=D_1(b-t_G)\ge0$;
for $b\in[t_0,t_B)$ she switches to Against only on $m=b$, gaining
$D_2[-(2\mu_B-1)-b]=D_2(t_B-b)\ge0$. Both pieces vanish at the band edges and
meet at $b=t_0$, where by (I1) each equals $2D_1(\mu_G-p)=2x(2Q-1)=h$.
\end{proof}

\begin{proposition}[Readership \statusP]\label{prop:readership}
A voter reads iff $\mathrm{VOI}(b_j)\ge c_r$. The reading set is the band
$[\ell,u]$ with $\ell=t_G+c_r/D_1$, $u=t_B-c_r/D_2$ (empty iff $h\le c_r$), and
its mass is
\begin{equation}\label{eq:A}
A(p;Q,c_r)\;=\;\frac{\bigl(2p(1-p)(2Q-1)-c_r\bigr)_+}{2a\,D_1D_2}
\;=\;\frac{\mu_G-\mu_B}{a}\Bigl(1-\frac{c_r}{h}\Bigr)_+ .
\end{equation}
$A$ is symmetric in $p$ about $\tfrac12$ and strictly increasing in $x=p(1-p)$
wherever positive; hence single-peaked at $p=\tfrac12$ and zero for consensus
proposals. Moreover, for the peak (ex-ante indifferent) voter, the net value of
reading is \emph{exactly}
\begin{equation}\label{eq:readrule}
V(\pi)\;=\;(1-\pi)\,I_H\;-\;c_r,\qquad I_H\equiv h(p,q)=2p(1-p)(2q-1),
\end{equation}
so she reads iff $\pi\le \bar\pi(p)\equiv 1-c_r/I_H(p)$.
\end{proposition}

\begin{proof}
The set $\{\mathrm{VOI}\ge c_r\}$ is the tent's cross-section at height $c_r$;
for a triangle of peak $h$, each side's width scales by $(1-c_r/h)$, giving mass
$(t_B-t_G)(1-c_r/h)_+/(2a)$; apply (I2), and (I1)+(I3) for the first form.
Monotonicity in $x$: differentiating \eqref{eq:A} with $D=Q(1-Q)+x(2Q-1)^2$,
\[
\frac{\partial A}{\partial x}
=\frac{2(2Q-1)\,Q(1-Q)+c_r(2Q-1)^2}{2a\,D^2}>0 .
\]
For \eqref{eq:readrule}: by Lemma~\ref{lem:tent} the peak voter's gross value is
$h(p,Q)=2x(2Q-1)$, and by \eqref{eq:Qeff}, $2Q-1=(1-\pi)(2q-1)$, so
$h(p,Q)=(1-\pi)h(p,q)=(1-\pi)I_H$.
\end{proof}

Equation \eqref{eq:readrule} \emph{is} the paper's equation
$V(\pi)=(1-\pi)I_H-c_r$ and its threshold $\pi^*=1-c_r/I_H$ --- previously
asserted for a representative delegate, now derived exactly for the marginal
(peak) reader, with the previously free parameter $I_H$ pinned to primitives:
$I_H=2p(1-p)(2q-1)$. Three consequences:
\begin{itemize}[itemsep=1pt]
\item \emph{The threshold is proposal-specific}: $\bar\pi(p)=1-c_r/[2p(1-p)(2q-1)]$.
Consensus proposals (low $x$) have low thresholds and stop being read first;
knife-edge proposals are read the longest. ``The forum dies from the tails
inward.''
\item \emph{Non-peak readers} have strictly smaller gross value (the tent), so
each voter has her own threshold below $\bar\pi(p)$; aggregate readership
\eqref{eq:A} declines in $\pi$ and reaches \emph{exactly zero} at the interior
contamination level $\bar\pi(p)<1$.
\item \emph{Readership falls with contamination}: with $c_r=0$,
$A=x(2q-1)(1-\pi)/(aD)$ --- near-linear in $(1-\pi)$ (exact numerator linearity;
the denominator's $\pi$-dependence is second order). This is the working notes'
new testable prediction; it survives the rebuild.
\end{itemize}

\section{The posting stage: volume, and why career concerns fail}
\label{sec:post}

\begin{proposition}[Posting volume \statusPA{} (Assumptions \ref{ass:truth}, \ref{ass:attention})]
\label{prop:volume}
Contributor $i$ posts iff $c_i\le R\,A(p)$, so equilibrium human posting mass is
\begin{equation}\label{eq:n}
n(p)=F\bigl(R\,A(p)\bigr),
\end{equation}
which inherits the shape of $A$: symmetric about $p=\tfrac12$, single-peaked
there, zero on consensus proposals with $h\le c_r$, and nonincreasing in
$d=|p-\tfrac12|$ (strictly decreasing wherever $0<RA<\bar c$). If $F$ is
near-linear at the origin and $c_r$ is small,
\[
n(p)\;\propto\;A(p)\;\approx\;\frac{2\,p(1-p)\,(2Q-1)}{a}\cdot\frac{1}{4D}
\;\propto\; p(1-p),
\]
up to the mild factor $\tfrac14 D^{-1}\in[1,\,(4Q(1-Q))^{-1}]$.
\end{proposition}

\begin{proof}
Immediate from Assumption~\ref{ass:attention} and
Proposition~\ref{prop:readership}. The approximation uses $D\to\tfrac14$ as
$p\to\tfrac12$ or $Q\to\tfrac12$.
\end{proof}

So ``posting volume $\propto p(1-p)$'' is delivered as a first-order
approximation to an exact closed form, with the exact form available for
structural work. Posting peaks on contested proposals \emph{because that is
where the audience is}, and the audience is there because that is where
information has value. One force (the tent peaking at $p=\tfrac12$) drives
reading, posting, and --- Section~\ref{sec:welfare} --- welfare.

\begin{remark}[Congestion variant]\label{rem:congestion}
If instead each of the $n$ posts receives only its \emph{share} of attention
(reward $R\,A/n$, readers spread over posts), the equilibrium condition becomes
$n=F(RA/n)$; with $F$ linear at the origin, $n\propto\sqrt{A}$. All shape
results below need only that $n$ is a strictly increasing transform of $A$, so
Theorem~\ref{thm:cov} is unaffected; only the exact locus ($p(1-p)$ vs.\ its
square root) changes --- a distinction the data cannot make anyway
(Section~\ref{sec:audit}).
\end{remark}

\subsection{The negative result: reputation for ability posts on the wrong proposals}

The fix roadmap proposed career concerns --- posting to attract delegated voting
power by signaling ability --- as the posting motive. We formalize it and prove
it fails, robustly.

Suppose contributors have ability $\alpha\in\{H,L\}$, $\Pro(H)=\gamma$, with
signal precisions $q_H>q_L>\tfrac12$; a contributor does not observe her own
type (the standard career-concerns case, Holmstr\"om 1999). Delegators observe a
poster's message and, ex post, $\theta$, and update; non-posters retain
reputation $\gamma$ (no inference from silence; see the remark below). The
poster's payoff is $\phi(\hat\gamma)-c_i$ where $\hat\gamma$ is her posterior
reputation and $\phi$ is \emph{any} strictly increasing function.

\begin{theorem}[Career concerns are wrong-signed \statusP]\label{thm:career}
Let $\bar q=\gamma q_H+(1-\gamma)q_L$. Then:
\begin{enumerate}[label=(\alph*),itemsep=1pt]
\item Posting leads to exactly two possible reputations, independent of $p$:
\begin{align*}
r_+&=\frac{\gamma q_H}{\gamma q_H+(1-\gamma)q_L}>\gamma
&&\text{(vindicated, } s=\theta\text{)},\\
r_-&=\frac{\gamma(1-q_H)}{\gamma(1-q_H)+(1-\gamma)(1-q_L)}<\gamma
&&\text{(wrong, } s\ne\theta\text{)}.
\end{align*}
\item The expected gain from posting signal $s$ is
\[
g(s,p)=\rho(s,p)\,\phi(r_+)+[1-\rho(s,p)]\,\phi(r_-)-\phi(\gamma),
\]
strictly increasing in the vindication probability
$\rho(s,p)=\Pro(\theta=s\mid s)$.
\item At $p=\tfrac12$, $\rho(s,\tfrac12)=\bar q$ for both signals; as
$p\to1$, $\rho(g,p)\to1$. Hence for every strictly increasing $\phi$, the
per-contributor posting gain at a foregone conclusion (majority signal) strictly
exceeds the gain at a toss-up, and total posting mass satisfies
\[
n^{\mathrm{cc}}(\tfrac12)\;<\;\lim_{p\to1}n^{\mathrm{cc}}(p).
\]
For linear $\phi$ the gain at $p=\tfrac12$ is exactly zero (so
$n^{\mathrm{cc}}(\tfrac12)=0$), and numerically $n^{\mathrm{cc}}(p)$ is V-shaped
in $p$ throughout the parameter range checked.
\end{enumerate}
Consequently career-concerns posting concentrates on \emph{consensus}
proposals, generating a nonnegative posts--margin covariance --- the opposite
of the data --- and cannot rationalize the paper's Proposition~1.
\end{theorem}

\begin{proof}
(a) Given $(s,\theta)$ the type-likelihood is $\Pro(s\mid\theta,\alpha)$, which
by symmetry of the precisions equals $q_\alpha$ if $s=\theta$ and $1-q_\alpha$
otherwise --- free of $p$. Bayes gives $r_\pm$; $r_+>\gamma>r_-$ follows from
$q_H>q_L$. (b) Immediate from (a). (c) $\rho(s,p)=\Pro(\theta=s\mid s)$ with the
mixture precision $\bar q$: at $p=\tfrac12$, $\rho=\bar q$ for both signals; as
$p\to1$, $\rho(g,p)=p\bar q/[p\bar q+(1-p)(1-\bar q)]\to1$ and the population
share of $g$-signals $\to1$. Since $\rho(g,\cdot)$ rises from $\bar q$ to $1$
and $\phi$ is strictly increasing, $g(g,p)\to\phi(r_+)-\phi(\gamma)>g(s,\tfrac12)$
for both $s$; posting mass is $\sum_s\Pro(s)F(g(s,p)_+)$, giving the strict
ranking. For linear $\phi$: by iterated expectations
$\E_\theta[\Pro(H\mid s,\theta)\mid s]=\Pro(H\mid s)$, and at $p=\tfrac12$ the
signal is uninformative about type ($\Pro(s\mid\alpha)=\tfrac12$ for both
$\alpha$), so $\Pro(H\mid s)=\gamma$ and $g=0$.
\end{proof}

The economics: on a toss-up your signal is (nearly) a coin flip \emph{whatever
your ability}, so being right is barely diagnostic in expectation --- and the
reputational payoffs $r_\pm$ from vindication are the same everywhere, while the
\emph{probability} of vindication is maximal when you hold the majority view on
a foregone conclusion. Career concerns therefore push posting toward safe
consensus, not contested debate. This kills the roadmap's step (ii) as stated,
and it also means the paper's delegate posting$\,\leftrightarrow\,$voting-power
correlation should \emph{not} be narrated as career-concerns evidence.

\begin{remark}[Robustness and caveats of the negative result]
The proof covers arbitrary increasing $\phi$ (so convex, tournament-style
delegation markets do not rescue the mechanism; numerically the V-shape persists
for $\phi(r)=r^3$ and $\phi=\sqrt{r}$, with the minimum near --- though for
convex $\phi$ not exactly at --- $p=\tfrac12$). Two modeling choices are
substantive: non-posters are not updated on (equilibrium inference from silence
is not modeled), and contributors do not know their own type. Relaxing either
changes the game but not the core force, which is (a): vindication payoffs are
$p$-invariant while vindication probability rises with consensus. We label the
endpoint comparison \textsc{proved} and the global V-shape
\textsc{numerically verified}.
\end{remark}

\section{The voting stage: interior margins, derived}
\label{sec:margin}

Now aggregate votes. Conditional on $\theta$, each reader's sampled message is
i.i.d.\ ($g$ with probability $\rho_G=Q$ if $\theta=G$; $\rho_B=1-Q$ if
$\theta=B$), so by Assumption~\ref{ass:lln} the For-share is deterministic given
$\theta$:
\[
\Phi(\theta)\;=\;\underbrace{\frac{a-u}{2a}}_{\text{non-readers with }b>u}
\;+\;\underbrace{A\cdot\rho_\theta}_{\text{readers voting For}},
\qquad\text{(non-readers with }b<\ell\text{ vote Against),}
\]
because within the reading band a $g$-message moves every reader to For and a
$b$-message to Against. The \emph{realized margin} is
$M(\theta)=|2\Phi(\theta)-1|$ and the expected margin is
$\E[M\mid p]=p\,M(G)+(1-p)\,M(B)$.

\begin{proposition}[Interior margins, closed form \statusP]\label{prop:margin}
Let $m_\theta\equiv2\Phi(\theta)-1$. In the active region ($h>c_r$),
\begin{align}
a\,D\;m_G&=Q(1-Q)(2p-1)\;+\;(2Q-1)^2x\;-\;c_r(1-p)(2Q-1),\label{eq:mG}\\
a\,D\;m_B&=Q(1-Q)(2p-1)\;-\;(2Q-1)^2x\;+\;c_r\,p\,(2Q-1),\label{eq:mB}
\end{align}
and when $h\le c_r$ (no readers) $m_G=m_B=(2p-1)/a$. For $p\ge\tfrac12$,
$m_G\ge0$ always; margins are interior (in $(0,1)$) for all interior $\mu$ by
Assumption~\ref{ass:interior}. Under $p\mapsto1-p$ the model is symmetric:
$m_B(1-p)=-m_G(p)$.
\end{proposition}

\begin{proof}
With band $[\ell,u]$, $\ell=t_G+c_r/D_1$, $u=t_B-c_r/D_2$: every reader votes
For on message $g$ and Against on message $b$ (the band lies inside
$[t_G,t_B]$), so $m_\theta=-u/a+2A\rho_\theta$. Substitute
$u=1-2\mu_B-c_r/D_2$ and \eqref{eq:A}, write
$2\mu_B-1=(\mu_G+\mu_B-1)-(\mu_G-\mu_B)$, and use (I2) together with three
one-line identities: $\mu_G+\mu_B-1=\mu_G-(1-\mu_B)=Q[pD_2-(1-p)D_1]/D$ with
$pD_2-(1-p)D_1=(1-Q)[p^2-(1-p)^2]=(1-Q)(2p-1)$, so
$\mu_G+\mu_B-1=Q(1-Q)(2p-1)/D$; and $D_1-Q=-(1-p)(2Q-1)$,
$D_1-(1-Q)=p(2Q-1)$ for the $c_r$ terms.
Sign of $m_G$ for $p\ge\tfrac12$: the negative term obeys
$c_r(1-p)(2Q-1)<2x(2Q-1)^2(1-p)\le(2Q-1)^2x$ using $c_r<h$ and $2(1-p)\le1$.
Symmetry: substitute $p\to1-p$ in \eqref{eq:mB} ($x$, $D$ invariant).
All formulas verified against direct numerical integration of the taste
distribution (300 random parameter draws, max abs.\ error $<4\times10^{-7}$).
\end{proof}

\begin{theorem}[Expected margin rises with consensus \statusP]\label{thm:margin}
$\E[M\mid p]$ is continuous, symmetric about $p=\tfrac12$, and \emph{strictly
increasing} in $d=|p-\tfrac12|$ on $(0,\tfrac12)$, for all
$Q\in(\tfrac12,1]$, $a\ge1$, $c_r\ge0$. Moreover
\begin{equation}\label{eq:floor}
\E[M\mid p]\;\ge\;\frac{|2p-1|}{a},
\end{equation}
with equality outside an interior contested window (in particular everywhere
once readership is zero), and $\E[M\mid\tfrac12]=(2Q-1)^2\bigl(1-\tfrac{2c_r}{2Q-1}\bigr)/a\ \ge 0$
at the center (with $c_r=0$: $(2Q-1)^2/a$).
\end{theorem}

\begin{proof}
Take $p\ge\tfrac12$ (symmetry). Write $\alpha=Q(1-Q)(2p-1)$,
$\beta=(2Q-1)^2x$, $\gamma_G=c_r(1-p)(2Q-1)$, $\gamma_B=c_rp(2Q-1)$, so
$aDm_G=\alpha+\beta-\gamma_G$, $aDm_B=\alpha-\beta+\gamma_B$. Three regimes as
$d$ rises from $0$:

\emph{(iii) Contested window} ($m_B<0$): here
$aD\,\E[M\mid p]=p(\alpha+\beta-\gamma_G)+(1-p)(\beta-\alpha-\gamma_B)
=\alpha(2p-1)+\beta-[p\gamma_G+(1-p)\gamma_B]$, and
$p\gamma_G+(1-p)\gamma_B=2c_r(2Q-1)x$, so with
$\kappa\equiv(2Q-1)^2/[Q(1-Q)]$ and
$\kappa_c\equiv(2Q-1)[(2Q-1)-2c_r]/[Q(1-Q)]\in(0,\kappa]$ (positive because
$2c_r<4x(2Q-1)\le 2Q-1$ in the active region),
\[
a\,\E[M\mid p]=\frac{4d^2+\kappa_c\,x}{1+\kappa\,x},\qquad x=\tfrac14-d^2 .
\]
Differentiating in $d^2$: the numerator of the derivative is
$(4-\kappa_c)(1+\kappa x)+\kappa(4d^2+\kappa_c x)=4+\kappa-\kappa_c>0$
(the $d^2$-terms cancel identically). Strictly increasing.

\emph{(ii) Consensus, active} ($m_B\ge0$, $h>c_r$): both margins positive, and
$aD[p\,m_G+(1-p)m_B]=\alpha+\beta(2p-1)+c_r(2Q-1)[-(1-p)p+(1-p)p]=(2p-1)D$, so
$\E[M\mid p]=(2p-1)/a$ \emph{exactly} --- strictly increasing.

\emph{(i) No readers} ($h\le c_r$): $\E[M\mid p]=(2p-1)/a$, strictly increasing.

Continuity: at the (iii)/(ii) boundary $m_B=0$ both expressions coincide; at
the (ii)/(i) boundary $\ell=u=t_0$ and $m_\theta\to(2p-1)/a$ smoothly. A
continuous function that is strictly increasing on each piece of a partition of
$[0,\tfrac12)$ into intervals is strictly increasing. The floor
\eqref{eq:floor}: in regimes (i)--(ii) it binds; in (iii),
$aD\,\E[M\mid p]-(2p-1)D=(1-2d)(\beta-\alpha)-2c_r(2Q-1)x
>(1-2d)\gamma_B-2c_r(2Q-1)x=c_r(2Q-1)[2p(1-p)-2x]=0$,
using $m_B<0\iff\beta-\alpha>\gamma_B$. Monotonicity was additionally verified
numerically on a $7\times3\times5$ parameter grid with $2\times10^4$ points of
$p$ each: zero violations.
\end{proof}

Objection (O1) is answered: margins are interior and continuously distributed
because voters disagree on \emph{values} ($a>0$); they are informative about
$p$ because the expected margin strictly increases in consensus $d$; and the
paper's shorthand ``margin $\propto|2p-1|$'' holds \emph{exactly} outside the
contested window (and as the leading term for weak signals, since
$\E[M\mid\tfrac12]=O((2Q-1)^2)$).

\section{The paper's Proposition 1, now a theorem}
\label{sec:prop1}

\begin{theorem}[Posts--margin covariance \statusP]\label{thm:cov}
Let proposals draw priors $p_i$ i.i.d.\ from any distribution such that
$d_i=|p_i-\tfrac12|$ is nondegenerate and readership $A(p_i)$ is nonconstant on
its support; let $\theta_i$ be drawn given $p_i$, and let
$(\text{Posts}_i,\text{Margin}_i)=(n(p_i),M_i)$ be the realized volume and
margin. Then
\[
\Cov(\text{Posts}_i,\ \text{Margin}_i)\;<\;0 .
\]
\end{theorem}

\begin{proof}
$n(p_i)$ is $\sigma(p_i)$-measurable, so by the law of total covariance
$\Cov(n,M)=\Cov\bigl(n(p),\E[M\mid p]\bigr)$ ($M$'s within-$p$ variation over
$\theta$ contributes nothing). By Propositions \ref{prop:readership} and
\ref{prop:volume}, $n$ is a nonincreasing function $\nu(d)$; by
Theorem~\ref{thm:margin}, $\E[M\mid p]$ is a strictly increasing function
$\rho(d)$. For i.i.d.\ copies $d,d'$,
$\Cov(\nu(d),\rho(d))=-\tfrac12\E[(\nu(d)-\nu(d'))(\rho(d)-\rho(d'))]$, where
every term in the expectation is $\le0$ (opposite monotonicities) and is
strictly negative with positive probability because $\nu$ is nonconstant and
$\rho$ is strictly increasing on a nondegenerate $d$.
\end{proof}

This is objection (O2) resolved on the paper's own terms: the negative
covariance between posting volume and margin of victory --- the paper's central
empirical object --- is a consequence of the equilibrium, with the posting
equilibrium sustained without any pivotality (Assumption~\ref{ass:attention})
and the margin process derived (Theorem~\ref{thm:margin}). Both conditional
moments are monotone transforms of the single latent $d_i$, which is what
licenses the paper's practice of reading margins as inverse measures of
contestedness.

\begin{remark}[The honest ceiling on empirical content]
Theorem~\ref{thm:cov} pins a \emph{sign}, not a locus. On the pre-break sample
($n=29$) the structural locus and a linear reduced form are statistically
indistinguishable (the working notes' $R^2$: $0.114$ vs $0.111$). The model's
discriminating content is cross-equation (the same $q$ appears in
\eqref{eq:A}, \eqref{eq:mG}--\eqref{eq:mB}, and the welfare window below), and
the working notes report those over-identifying tests as underpowered on this
sample. Deriving the model buys internal rigor, not additional empirical
verification --- this should be said in the paper.
\end{remark}

\section{Welfare: the exact value of the forum, and the fate of $4p(1-p)(q-\frac12)^2$}
\label{sec:welfare}

The DAO's decision payoff is the probability the outcome matches $\theta$
(pass iff $G$). Without a forum, the vote aggregates only tastes and the prior:
the outcome follows the prior action, correct with probability $\max(p,1-p)$.

\begin{proposition}[Exact decision value of the forum \statusP]\label{prop:welfare}
Define the \emph{flip window}
\[
\Omega(Q,c_r)\;=\;\Bigl\{p:\ (2Q-1)^2p(1-p)\;-\;c_r(2Q-1)\max(p,1-p)\;>\;Q(1-Q)\,|2p-1|\Bigr\},
\]
a symmetric interval around $p=\tfrac12$, nonempty iff $c_r<\tfrac{2Q-1}{2}$,
expanding to $(0,1)$ as $Q\to1,c_r\to0$ and shrinking to $\varnothing$ as
$Q\to\tfrac12$. In equilibrium:
\begin{enumerate}[label=(\alph*),itemsep=1pt]
\item For $p\in\Omega$: $m_G>0>m_B$ --- the vote passes good proposals and
rejects bad ones \emph{with probability one} (full information aggregation, the
Feddersen--Pesendorfer property delivered by the LLN across readers).
\item For $p\notin\Omega$: the outcome coincides with the prior action.
\item Hence the forum's gross decision value is
\begin{equation}\label{eq:V}
V(p;Q,c_r)\;=\;\min(p,\,1-p)\cdot\ind{p\in\Omega(Q,c_r)},
\end{equation}
and net welfare subtracts participation costs
$W=\Pro(\mathrm{correct})-c_r A(p)-\int_0^{RA(p)}c\,dF(c)$.
\end{enumerate}
\end{proposition}

\begin{proof}
For $p\ge\tfrac12$ in the active region, $m_G>0$ always
(Proposition~\ref{prop:margin}), so correctness turns on $m_B<0$, i.e.
$\beta-\gamma_B>\alpha$, which is the displayed window condition with
$\max(p,1-p)=p$; the $p\le\tfrac12$ branch is the mirror condition $m_G>0$.
Window $\subset$ active region: $m_B<0$ requires $\beta>\gamma_B$, i.e.
$x(2Q-1)>c_rp$, which for $p\ge\tfrac12$ implies $h=2x(2Q-1)>2c_rp\ge c_r$.
Nonemptiness iff the condition holds at $p=\tfrac12$:
$(2Q-1)^2/4-c_r(2Q-1)/2>0\iff c_r<(2Q-1)/2=h(\tfrac12,Q)$, which is also
exactly the condition that \emph{any} proposal has readers. (b): outside
$\Omega$ but active, $m_G$ and $m_B$ share the prior's sign, so the outcome is
the prior action; inactive is immediate. (c): conditional correctness is
$p\cdot1+(1-p)\cdot1=1$ inside and $\max(p,1-p)$ outside; subtract the prior
baseline.
\end{proof}

\paragraph{Verdict on the paper's formula.} The kernel
$4p(1-p)(q-\tfrac12)^2$ is \emph{not} the value of the forum in any version of
the model:
\begin{itemize}[itemsep=1pt]
\item In the derived two-tier model the exact object is \eqref{eq:V}: a tent
$\min(p,1-p)$ truncated to an interior window --- \emph{closeness enters as the
level, informativeness as the width}. AI dilution ($Q\to\tfrac12$ via
\eqref{eq:Qeff}) destroys value by \emph{narrowing the window} until it
vanishes at $\bar\pi(\tfrac12)$; the collapse destroys $\min(p,1-p)$ of decision
quality per proposal, maximal ($=\tfrac12$) on contested proposals. The
qualitative claim the paper builds on --- AI destroys the most welfare where
governance is most contested --- survives exactly.
\item For a \emph{single} common signal of precision $q$ (closest to the old
representative-delegate framing), the exact value is the Blackwell tent
$(q-\max(p,1-p))_+$, and the exact belief-accuracy (variance-reduction) value is
$[p(1-p)(2q-1)]^2/(D_1D_2)$, which is $16[p(1-p)]^2(q-\tfrac12)^2$ to leading
order --- note the \emph{square} on $p(1-p)$. The paper's monomial matches
neither; it is a smooth mnemonic with the right single-peakedness and the right
$(q-\tfrac12)^2$ scaling but no decision-theoretic pedigree.
\end{itemize}
\emph{Recommendation:} replace the formula with \eqref{eq:V} (decision-quality
story, matches the rest of the model), or --- if a smooth kernel is wanted for
prose --- present $4p(1-p)(q-\tfrac12)^2$ explicitly as a heuristic
interpolation, never as a result. Since the words-based calibration never used
this formula (it runs off $I_H$, $c_r$, $\pi^*$ --- see
Section~\ref{sec:calibration}), replacing it costs nothing quantitative.

\begin{remark}[Why full aggregation inside the window is the right idealization]
Statement (a) is the continuum idealization of a Condorcet effect: finitely many
readers make the probability of a correct outcome $1-O(e^{-cN})$ inside the
window. The knife-edge flavor of ``probability one'' is Assumption~\ref{ass:lln},
not deep structure; what is robust is that decision quality is high inside the
window and prior-driven outside.
\end{remark}

\section{The AI shock: dilution, margin compression, and the reading collapse}
\label{sec:shock}

AI enters through two channels, both now derived: the sampled-post precision
falls (equation \eqref{eq:Qeff}), and the reading calculus \eqref{eq:readrule}
deteriorates.

\begin{proposition}[Comparative statics in contamination \statusP]\label{prop:shock}
As $\pi$ rises (equivalently $Q\downarrow\tfrac12$):
\begin{enumerate}[label=(\alph*),itemsep=1pt]
\item \emph{Readership} $A(p;\pi)$ falls, with exactly linear numerator
$\bigl(2x(2q-1)(1-\pi)-c_r\bigr)_+$, reaching zero at
$\bar\pi(p)=1-c_r/I_H(p)<1$: an \emph{interior} contamination level at which
the forum's audience is exactly zero. Consensus proposals hit their thresholds
first.
\item \emph{Posting} $n(p;\pi)=F(RA)$ falls with readership --- the audience
channel: AI does not need to out-argue humans, it only needs to thin the
readership that pays for human effort.
\item \emph{Margins} compress toward the no-forum benchmark: $\E[M\mid p]$ is
nondecreasing in $Q$ pointwise, falling to the floor $|2p-1|/a$; the contested
window $\Omega$ narrows to $\varnothing$.
\item \emph{Welfare}: $V(p;Q_{\mathrm{eff}},c_r)$ collapses window-first; at
$\pi\ge\bar\pi(\tfrac12)$ the forum's entire decision value is gone even though
$1-\pi$ of posts are genuine.
\end{enumerate}
\end{proposition}

\begin{proof}
(a) is Proposition~\ref{prop:readership} with \eqref{eq:Qeff}; (b) immediate.
(c) The floor is attained: $\E[M\mid p]\ge|2p-1|/a$ always
(Theorem~\ref{thm:margin}), with equality everywhere once $h\le c_r$, which
occurs for every $p$ when $Q$ is close enough to $\tfrac12$; so as
$\pi\to\bar\pi(\tfrac12)$ the entire margin process converges to the no-forum
benchmark. Pointwise monotonicity in $Q$: for $c_r=0$,
$a\E[M|p]=(4d^2+\kappa x)/(1+\kappa x)$ in the window and
$\partial(a\E)/\partial\kappa=x(1-4d^2)/(1+\kappa x)^2\ge0$ with
$\kappa=(2Q-1)^2/[Q(1-Q)]$ increasing in $Q$ --- proved; for $c_r>0$ the two
channels ($\kappa_c\uparrow$, $\kappa\uparrow$) push in opposite directions and
we verify $\partial\E[M|p]/\partial Q\ge0$ numerically (grid over
$a\in\{1,1.5,3\}$, $c_r\le0.2$, $p\in[0.5,0.99]$, $2{,}000$ values of $Q$:
zero violations). (d) The window condition is monotone in $Q$ and $c_r$.
\end{proof}

Part (a) is the precise form the paper's Proposition~2 (``unverifiability
threshold'') takes in the rebuilt model, and the honest statement is slightly
different from the draft's:
\begin{itemize}[itemsep=1pt]
\item \emph{At the individual level} the paper's discreteness survives exactly:
each reader's problem is binary, and the marginal reader's rule is
\eqref{eq:readrule} with a hard threshold.
\item \emph{At the aggregate level}, readership declines continuously and hits
\emph{exactly zero at an interior} $\bar\pi<1$ --- a collapse at finite
contamination, but a kinked continuous decline rather than a jump. The
genuinely \emph{discontinuous} collapse the abstract advertises is restored ---
and made robust --- once contamination is endogenous:
Theorem~\ref{thm:tipping}.
\item The draft's distinction between expected contamination $\pi$ (the object
in the reading rule) and realized share $\lambda$ carries over verbatim: the
decision-relevant $\pi$ is the reader's forward-looking expectation about the
post she is about to read, which jumps at a capability release ahead of
$\lambda$. This interpretation layer is unchanged by the rebuild and remains the
mechanism for collapse at $\lambda\approx8\%$.
\end{itemize}

\section{Endogenous contamination: a tipping theorem}
\label{sec:tipping}

Close the loop: contamination depends on how much humans post, which depends on
the audience, which depends on contamination. Fix a proposal type (contestedness
$x$; think of the representative contested proposal, $x=\tfrac14$) and let AI
post mass be $z\ge0$. An equilibrium is a triple $(A,n,\pi)$ with
\[
n=F(RA),\qquad \pi=\frac{z}{z+n},\qquad
A=\mathcal{A}(\pi)\equiv\frac{\bigl(2x(2q-1)(1-\pi)-c_r\bigr)_+}{2a\,D(\pi)},
\]
$D(\pi)=Q(1-Q)+x(2Q-1)^2$ at $Q=\tfrac12+(1-\pi)(q-\tfrac12)$. Equivalently, a
fixed point of the audience map
$T_z(A)\equiv\mathcal{A}\bigl(z/(z+F(RA))\bigr)$.

\begin{theorem}[Discontinuous, hysteretic collapse \statusPA{} (representative proposal; exogenous $z$; Assumptions \ref{ass:lln}--\ref{ass:viable})]
\label{thm:tipping}
Let $c_r\in(0,I_H)$ and $\bar\pi=1-c_r/I_H$. Then:
\begin{enumerate}[label=(\alph*),itemsep=1pt]
\item \emph{Collapse always coexists.} $(A,n,\pi)=(0,0,1)$ is an equilibrium for
every $z>0$.
\item \emph{Dead zone.} $T_z(A)=0$ for all $A\le A_z\equiv F^{-1}\bigl(\min\{z\tfrac{1-\bar\pi}{\bar\pi},\,1\}\bigr)/R$,
and $A_z>0$ for every $z>0$: audiences below $A_z$ cannot sustain enough human
posting to keep expected reading value positive.
\item \emph{Tipping point.} $T_z$ is nondecreasing in $A$ and strictly
decreasing in $z$, so
$z^*\equiv\sup\{z:\text{a healthy equilibrium }A>0\text{ exists}\}$
is well defined, with $0<z^*\le\frac{\bar\pi}{1-\bar\pi}\,n_{\max}$
($n_{\max}=F(\bar c)$): once AI post mass exceeds $\frac{\bar\pi}{1-\bar\pi}$
times maximal human volume, only collapse survives. For every $z\in(0,z^*)$ the
healthy and collapsed equilibria \emph{coexist} (bistability); for $z>z^*$
collapse is unique.
\item \emph{Discontinuity.} Every healthy equilibrium satisfies $A>A_z$; hence
as $z$ crosses $z^*$ the audience jumps from at least $A_{z^*}>0$ to $0$ ---
never a continuous fade-out.
\item \emph{Hysteresis.} Under best-response dynamics $A_{t+1}=T_z(A_t)$, the
collapsed state is locally absorbing for every $z>0$ (by (b)); after a collapse,
reducing $z$ --- even nearly to zero --- does not restore the forum. Recovery
requires coordinated re-entry of audience mass exceeding $A_z$, not merely
undoing the shock.
\item \emph{Collapse below the reading threshold.} At any healthy equilibrium,
realized contamination satisfies $\pi=z/(z+n)<\bar\pi$ strictly. The forum
therefore always tips at a realized AI share strictly below the level
$\bar\pi$ at which reading value would mechanically reach zero.
\end{enumerate}
\end{theorem}

\begin{proof}
(a) $A=0\Rightarrow n=F(0)=0\Rightarrow\pi=1\Rightarrow Q=\tfrac12\Rightarrow
\mathcal{A}=0$. (b) $\mathcal{A}(\pi)>0$ iff $\pi<\bar\pi$ iff
$F(RA)>z(1-\bar\pi)/\bar\pi$. (c) Monotonicity of $T_z$ in $A$: the numerator of
$\mathcal{A}$ is increasing in $1-\pi$ and $\pi$ is decreasing in $A$; the
denominator $D$ is \emph{decreasing} in $Q$ since
$\partial D/\partial Q=2(2Q-1)(4x-1)\cdot\tfrac12\le0$ for $x\le\tfrac14$;
strict decrease in $z$ is immediate. Finiteness of $z^*$: a healthy equilibrium
needs $F(RA)>z(1-\bar\pi)/\bar\pi$ with $F\le n_{\max}$. Positivity: at $z=0$
the healthy audience is $A_0=\mathcal{A}(0)>0$ (Assumption~\ref{ass:viable});
for any $A'<A_0$, $T_z(A')\to\mathcal{A}(0)>A'$ as $z\to0$, and a nondecreasing
self-map of $[A',A_{\max}]$ has a fixed point (Tarski); the set of $z$ with a
healthy equilibrium is a down-set by pointwise monotonicity in $z$.
Bistability on $(0,z^*)$: (a) plus the down-set property. (d) By (b), any
nonzero fixed point exceeds $A_z$, and $A_z$ is nondecreasing in $z$.
(e) $T_z\equiv0$ on $[0,A_z]$, $A_z>0$. (f) $A=T_z(A)>0$ requires
$\pi(A)<\bar\pi$ by (b).
\end{proof}

\begin{remark}[What this strengthens, and what it does not deliver]
The working notes' Proposition~3 obtained tipping and hysteresis \emph{only
under an S-shaped adoption curve} and flagged that with concave adoption the
decline is smooth. Theorem~\ref{thm:tipping} removes that condition: a positive
reading cost alone creates the dead zone that makes collapse discontinuous,
locally absorbing, and coexisting with health for every $z\in(0,z^*)$ --- for
\emph{any} continuous strictly increasing $F$. The reading cost that generates
the unverifiability threshold is the same primitive that makes the collapse
irreversible; that unification is new and, we think, worth a proposition in the
paper (or the lead result of a spin-off). Numerical benchmark ($x=\tfrac14$,
$q=0.75$, $a=1.2$, uniform $F$ on $[0,0.25]$, $R=1$): $c_r=0.02$ gives
$\bar\pi=0.92$ but collapse at $z^*\approx0.86$ with realized contamination
$\pi\approx0.69$ and an audience jump $\approx0.10\to0$; raising $c_r$ to
$0.08$ moves collapse to realized $\pi\approx0.41$. Thin reading margins mean
collapse strikes far below saturation. What the theorem does \emph{not} deliver:
the observed $\lambda\approx8\%$ is still lower than these equilibrium
contamination levels under the benchmark parameters --- matching the \emph{date}
of the collapse at 8\% realized share still requires the expectations channel
($\pi$ jumping ahead of $\lambda$ at the capability release), exactly as the
paper currently argues. The theorem's contribution is that collapse at
low-and-partial contamination is an equilibrium property, not a fragility of
one calibration. Caveats: representative-proposal reduction (aggregation across
a distribution of $p$ is not done); $z$ exogenous; stability/hysteresis defined
for best-response dynamics.
\end{remark}

\section{Calibration audit: what survives of $\pi^*\ge0.082$ and 858 words}
\label{sec:calibration}

The paper's Section 5.5 runs entirely off the reading rule, so we audit each
step against the derived model.

\textbf{Step 1: the threshold identity.} $\pi^*=1-c_r/I_H$, hence
``$\pi^*$ equals the net-surplus share of a post.'' \emph{Survives exactly}
(equation \eqref{eq:readrule}), now with $I_H=2p(1-p)(2q-1)$ pinned by
primitives and the representative delegate mapped to the peak (ex-ante
indifferent) reader on the representative contested proposal. \statusP

\textbf{Step 2: the bound $\pi^*\ge\lambda_{\text{break}}=0.082$.} The logic
(reading was still happening at the break $\Rightarrow\pi\le\pi^*$; expectations
run ahead of a rising realized share $\Rightarrow\pi\ge\lambda$) involves no
structure beyond the threshold rule and is unchanged. \emph{Survives.}
Theorem~\ref{thm:tipping}(f) adds an independent reason collapse precedes
saturation. \statusPA{} (the $\pi\ge\lambda$ step is an assumption about
expectations, as the paper already flags)

\textbf{Step 3: per-post surplus $I_H-c_r\ge\pi^*I_H\ge0.082\times237\approx19.4$
words.} Survives for the marginal reading decision of the peak reader with the
paper's normalization $I_H=c_p=237$ words. \statusPA{} (normalization)

\textbf{Step 4: aggregation $\times\,\bar n=44.2$ posts to
$\ge858$ words/delegate/proposal, and $\times K=62$ to $\approx43{,}000$.}
\emph{This is where the derived model withholds full support.} In the one-post
sampling protocol a reader consumes one post, and the derived aggregate surplus
is far smaller. If delegates read whole threads (the paper's implicit protocol),
per-post surplus is constant at $19.4$ only if successive posts do not become
redundant. With conditionally i.i.d.\ signals, marginal reading value
\emph{declines} in posts already read, and the sum is strictly less than
$44.2\times19.4$; only the first-post bound ($\ge19.4$ words per delegate per
proposal) is fully derived. The $858$ figure survives under:

\begin{assumption}[Non-redundancy, NR]\label{ass:nr}
Posts address (approximately) distinct payoff-relevant aspects of a proposal,
so the informational value of each successive post is undiminished and the
threshold calculus \eqref{eq:readrule} applies post-by-post.
\end{assumption}

Under NR the paper's arithmetic goes through unchanged. NR is defensible for
governance threads (posts flag distinct risks, precedents, technical issues)
but it is an assumption, not a theorem; a redundant-signal version of the same
calculation would justify only the $19.4$-word floor plus unquantified
inframarginal surplus. Additionally, averaging over the reading band rather
than the peak reader halves per-reader surplus (tent geometry). \emph{Verdict:
858 and 43{,}000 are \statusPA{} (NR + whole-thread protocol + peak-reader
mapping) --- i.e., exactly the ``illustrative calibration under stated
assumptions'' framing of agenda Decision~2, which this analysis independently
recommends. The units, the bound logic, and the 19.4 floor are solid; the
headline magnitudes are assumption-carrying.}

\begin{remark}[Multi-post reading \statusS]\label{rem:kpost}
A full sequential-reading extension (read posts one at a time at $c_r$ each,
stop optimally) would endogenize how many of the $\bar n$ posts are read and
let the data separate NR from redundancy via the observed reads-per-post decay
within threads. This is the natural next modeling increment if the calibration
is to be upgraded from illustrative to identified; it is not done here.
\end{remark}

\section{Model--data audit}
\label{sec:audit}

Where the empirics constrain the theory, does the derived model deliver?

\begin{center}
\begin{tabular}{p{5.6cm}p{8.6cm}}
\toprule
\textbf{Empirical shape} & \textbf{Delivered?}\\
\midrule
Posting volume peaked on contested proposals ($\propto p(1-p)$) &
Yes: exact closed form \eqref{eq:A}--\eqref{eq:n}, $\propto p(1-p)$ to first
order. Locus not separable from linear reduced form at $n=29$ --- claim the
sign, not the shape.\\[3pt]
Margin increasing in consensus ($\propto|2p-1|$) &
Yes: Theorem~\ref{thm:margin}; exact floor $|2p-1|/a$, strict monotonicity
everywhere.\\[3pt]
$\Cov(\text{posts},\text{margin})<0$ pre-AI; $\to0$ post &
Yes: Theorem~\ref{thm:cov}; post-collapse both $n$ and the $p$-sensitivity of
$M$ shut down, so the covariance dies.\\[3pt]
Collapse at low realized AI share &
Qualitatively yes: interior threshold $\bar\pi<1$; equilibrium tipping strictly
below $\bar\pi$ (Theorem~\ref{thm:tipping}f); the specific 8\% still needs the
expectations channel $\pi>\lambda$.\\[3pt]
Reads per post decline post-shock &
Yes, new derived prediction: $A$ near-linear in $(1-\pi)$. Empirically right
sign, underpowered (working notes \S6.1).\\[3pt]
Delegates using AI hold less voting power &
Outside the formal model. Consistent with the attention channel (uninformative
posts attract fewer readers/delegators) --- narrative only. The career-concerns
reading is \emph{ruled out} by Theorem~\ref{thm:career}.\\
\midrule
\multicolumn{2}{l}{\textbf{Honest misses (must be acknowledged in the paper):}}\\
\midrule
Post-break margins \emph{rose} (0.72$\to$0.91) &
Model predicts compression toward $|2p-1|/a$ on contested proposals, not a
rise. A rise requires a composition shift in the $p_i$ distribution (fewer
contested proposals reaching votes post-shock) --- plausible, untested, outside
the model.\\[3pt]
Humans still write $\approx$8.9k words/proposal post-break &
Model's $n\to0$ at collapse overshoots. Patch: an audience-independent posting
motive $\varepsilon>0$ (intrinsic/expressive) leaves all propositions intact
and delivers residual volume; not formalized here.\\
\bottomrule
\end{tabular}
\end{center}

\section{Verdict for the Thursday meeting}
\label{sec:verdict}

\subsection*{Fully proved (from primitives, given the model's stated assumptions)}
\begin{itemize}[itemsep=1pt]
\item Interior margins with closed forms; expected margin strictly increasing in
consensus, floor $|2p-1|/a$ (Prop.~\ref{prop:margin}, Thm.~\ref{thm:margin}).
\item Readership tent and closed form; the paper's reading rule
$V(\pi)=(1-\pi)I_H-c_r$ and threshold $\pi^*=1-c_r/I_H$, with
$I_H=2p(1-p)(2q-1)$ (Prop.~\ref{prop:readership}).
\item The paper's Proposition~1, $\Cov(\text{posts},\text{margin})<0$
(Thm.~\ref{thm:cov}).
\item Career concerns as posting motive are wrong-signed for any increasing
reputation payoff --- the roadmap's step (ii) is refuted and replaced
(Thm.~\ref{thm:career}).
\item Exact welfare: forum value $=\min(p,1-p)$ on an interior window
(Prop.~\ref{prop:welfare}). The old kernel $4p(1-p)(q-\tfrac12)^2$ is
\emph{not derivable} and must be replaced or demoted to a heuristic.
\item AI comparative statics: dilution equation from the sampling protocol;
window-first welfare collapse; margin compression; readership near-linear in
$(1-\pi)$ (Prop.~\ref{prop:shock}).
\end{itemize}

\subsection*{Proved under stated assumptions}
\begin{itemize}[itemsep=1pt]
\item The whole edifice rests on Assumptions \ref{ass:sincere} (decision-theoretic
voting/reading) and \ref{ass:attention} (attention rewards): these \emph{bypass}
pivotality rather than solve it. They are standard, defensible, and must be
stated in the paper as assumptions with their microfoundations discussed ---
this is the honest answer to the referee, and it is a good one: given two
behavioral primitives, everything else is derived.
\item Truthful posting (Assumption~\ref{ass:truth}; reputational cheap-talk
justification, not modeled).
\item Tipping Theorem~\ref{thm:tipping}: discontinuous, hysteretic,
below-threshold collapse --- proved, under a representative-proposal reduction
with exogenous AI mass and best-response dynamics. Strictly stronger than the
S-shape-conditional sketch in the working notes.
\item Calibration: threshold identity and $\pi^*\ge0.082$ survive; the 19.4
words/post surplus floor survives; \textbf{858 words/delegate and 43{,}000
words/proposal survive only under Assumption NR} (non-redundant posts +
whole-thread reading). Present as ``illustrative under stated assumptions''
(confirms agenda Decision~2's recommendation).
\end{itemize}

\subsection*{Open / sketch}
\begin{itemize}[itemsep=1pt]
\item Multi-post sequential reading (would upgrade the calibration; Remark~\ref{rem:kpost}).
\item Aggregating the tipping theorem across a cross-section of proposals;
endogenous AI supply $z$; equilibrium inference from silence in the
career-concerns game.
\item The post-break \emph{rise} in margins and residual human posting: need a
composition story and an $\varepsilon$-posting motive respectively; neither is
formalized.
\end{itemize}

\subsection*{Recommendation}
``Derive fully'' is \textbf{viable for the core} and no longer risky: the
derivations above are complete, checked numerically end-to-end, and produce the
paper's Propositions 1 and 2 as theorems. The honest package for the paper is a
hybrid:
\begin{enumerate}[itemsep=1pt]
\item \emph{Replace \S5.1--5.3 with the derived two-tier model} (Sections
\ref{sec:model}--\ref{sec:prop1}, \ref{sec:shock} here), stating Assumptions
\ref{ass:sincere} and \ref{ass:attention} candidly. Est.\ 1--1.5 weeks of
writing; no new mathematical risk remains.
\item \emph{Replace the welfare formula} with Proposition~\ref{prop:welfare}
(or demote the monomial to an explicit heuristic). Mandatory --- the old formula
is now known (to us) to be underivable, and a referee can find this.
\item \emph{Reframe the calibration} as illustrative under stated assumptions,
keeping the threshold identity and the $0.082$ bound as the load-bearing
numbers and stating NR when quoting 858/43{,}000. (= agenda Decision~2.)
\item \emph{Tipping theorem}: include compactly as the mechanism capstone
(it is now a theorem, not a possibility sketch) \emph{or} hold it for the
theory follow-up; either is defensible. If included, keep the
expectations-vs-realization distinction as the bridge to the 8\% fact.
\item \emph{Do not} use career-concerns language for the posting motive or the
delegate cross-section anywhere in the paper; Theorem~\ref{thm:career} makes
that framing indefensible.
\end{enumerate}
Estimated total: 2--3 weeks to a referee-solid Section 5, most of it exposition
rather than mathematics. The alternative (``reframe as illustrative'') is no
longer the safer option: the derivation exists, is checkable, and the
illustrative framing would leave the strongest referee line unanswered while
this note shows it answerable.

\begin{thebibliography}{9}
\bibitem{akerlof1970} Akerlof, G. (1970). The market for ``lemons'': Quality
uncertainty and the market mechanism. \emph{Quarterly Journal of Economics},
84(3), 488--500.
\bibitem{dewatripont1999} Dewatripont, M., and J. Tirole (1999). Advocates.
\emph{Journal of Political Economy}, 107(1), 1--39.
\bibitem{fp1997} Feddersen, T., and W. Pesendorfer (1997). Voting behavior and
information aggregation in elections with private information.
\emph{Econometrica}, 65(5), 1029--1058.
\bibitem{fs2006} Feddersen, T., and A. Sandroni (2006). A theory of
participation in elections. \emph{American Economic Review}, 96(4), 1271--1282.
\bibitem{gs1980} Grossman, S., and J. Stiglitz (1980). On the impossibility of
informationally efficient markets. \emph{American Economic Review}, 70(3),
393--408.
\bibitem{holmstrom1999} Holmstr\"om, B. (1999). Managerial incentive problems:
A dynamic perspective. \emph{Review of Economic Studies}, 66(1), 169--182.
\bibitem{judd1985} Judd, K. (1985). The law of large numbers with a continuum
of IID random variables. \emph{Journal of Economic Theory}, 35(1), 19--25.
\bibitem{lohmann1994} Lohmann, S. (1994). Information aggregation through
costly political action. \emph{American Economic Review}, 84(3), 518--530.
\bibitem{os2006} Ottaviani, M., and P. S{\o}rensen (2006). Reputational cheap
talk. \emph{RAND Journal of Economics}, 37(1), 155--175.
\end{thebibliography}

\end{document}
